\documentclass[11pt,a4paper]{article}
\usepackage[utf8]{inputenc}
\usepackage[spanish]{babel}
\usepackage{geometry}
\geometry{top=2.5cm,bottom=2.5cm,left=2.5cm,right=2.5cm}
\usepackage{fancyhdr}
\usepackage{amsmath,amssymb}
\usepackage{booktabs}
\usepackage{xcolor}

\definecolor{navyblue}{RGB}{31,78,121}
\definecolor{darkslate}{RGB}{47,85,151}

\pagestyle{fancy}
\fancyhf{}
\rhead{\small\textcolor{darkslate}{MCH Río Blanco 100 kW - Puno}}
\lhead{\small\textcolor{darkslate}{Expediente Técnico de Ingeniería}}
\rfoot{\small Página \thepage}
\lfoot{\small 2024}

\begin{document}

\begin{center}
    {\LARGE\bfseries\textcolor{navyblue}{ANEXO B: REGISTROS HIDROMETEOROLÓGICOS Y TABLAS ESTADÍSTICAS}}\\[0.6em]
    {\large\bfseries Series Históricas Multianuales SENAMHI y Aforos de Caudales en el Río Blanco}\\[1.2em]
    \rule{\linewidth}{1.2pt}
\end{center}


\section*{1. Estadísticos Climatológicos Multianuales (Estación Ollachea / San Gabán)}
\begin{table}[h!]
\centering
\small
\begin{tabular}{lcccccc}
\toprule
\textbf{Mes} & \textbf{P. Media [mm]} & \textbf{P. Máx 24h [mm]} & \textbf{T. Media [$^\circ$C]} & \textbf{T. Mín [$^\circ$C]} & \textbf{T. Máx [$^\circ$C]} & \textbf{HR [\%]} \\
\midrule
Enero & 262.4 & 68.5 & 16.2 & 9.5 & 22.8 & 84.2 \\
Febrero & 248.1 & 74.2 & 15.8 & 9.8 & 22.4 & 85.6 \\
Marzo & 212.5 & 62.0 & 16.0 & 9.2 & 22.7 & 83.1 \\
Abril & 94.8 & 38.4 & 15.2 & 7.6 & 23.1 & 76.5 \\
Mayo & 44.2 & 22.0 & 14.1 & 5.1 & 23.4 & 68.2 \\
Junio & 24.6 & 16.5 & 13.2 & 3.4 & 23.0 & 62.1 \\
Julio & 21.0 & 15.0 & 12.8 & 2.8 & 22.9 & 59.8 \\
Agosto & 36.5 & 20.4 & 13.9 & 4.2 & 23.5 & 63.4 \\
Septiembre & 66.8 & 28.5 & 15.0 & 6.5 & 24.2 & 67.9 \\
Octubre & 112.4 & 42.0 & 16.1 & 8.0 & 24.8 & 72.4 \\
Noviembre & 162.0 & 51.2 & 16.4 & 8.8 & 24.5 & 76.8 \\
Diciembre & 234.7 & 64.8 & 16.3 & 9.4 & 23.6 & 82.0 \\
\midrule
\textbf{Anual} & \textbf{1520.0} & \textbf{74.2} & \textbf{15.1} & \textbf{2.8} & \textbf{24.8} & \textbf{73.5} \\
\bottomrule
\end{tabular}
\end{table}

\section*{2. Régimen de Caudales Clasificados (Curva de Persistencia)}
\begin{table}[h!]
\centering
\small
\begin{tabular}{lcccc}
\toprule
\textbf{Percentil} & \textbf{Caudal Río [m$^3$/s]} & \textbf{Q Ecológico [m$^3$/s]} & \textbf{Q Turbinado [m$^3$/s]} & \textbf{Potencia Neta [kW]} \\
\midrule
Q10\% & 1.094 & 0.030 & 0.125 & 104.42 \\
Q25\% & 0.748 & 0.030 & 0.125 & 104.42 \\
Q50\% & 0.333 & 0.030 & 0.125 & 104.42 \\
Q75\% & 0.210 & 0.030 & 0.125 & 104.42 \\
Q78\% (Diseño) & 0.198 & 0.030 & 0.125 & 104.42 \\
Q85\% & 0.173 & 0.030 & 0.112 & 93.50 \\
Q90\% & 0.152 & 0.030 & 0.105 & 87.70 \\
Q95\% & 0.126 & 0.030 & 0.096 & 80.20 \\
\bottomrule
\end{tabular}
\end{table}


\end{document}
